\section{Introduction}
\IEEEPARstart{A}{chieving} human-level dexterity for robots has long been a central challenge in robotic research. The prospect of bimanual robotic systems with dexterous hands that mirror human sensorimotor dexterity holds the promise of seamlessly integrating robots into daily human activities and environments. Despite notable progress, how to capitalize on the abundance of available human data and develop algorithms suited to intricate and high-precision dexterous manipulation remains underexplored. 


Humans continuously generate diverse bimanual manipulation data, inherently serving as natural guidance for robots to emulate human-like behaviors. Several previous studies~\cite{zhou2025you,kim2025uniskill,lum2025crossing,dan2025x,wangmimicplay} have attempted to extract trajectories of human hands and manipulated objects from video data, subsequently applying them to robotic manipulation tasks. Nevertheless, these methods have predominantly targeted robots equipped with simple parallel gripper end effectors, failing to generalize effectively to dexterous hands due to the vastly greater complexity of action space. Despite recent advances that utilize kinematic retargeting approaches to produce human-like robotic motions~\cite{qiu2025humanoid,qin2023anyteleop,yang2025egovla,shaw2024learning,shaw2023videodex}, these approaches still fall short in achieving physically-aware pose retargeting and bridging the embodiment gap to derive feasible robot actions capable of successfully accomplishing the intended tasks. A critical limitation lies in omitting the modeling of interactions between robotic hands and manipulated objects, a fundamental component of manipulation tasks. Consequently, neglecting these interactions undermines the robot's ability to fully understand and adapt to the dynamics of manipulation scenarios.
 

Therefore, in an attempt to address the aforementioned challenges, recent approaches have begun leveraging Reinforcement Learning~(RL) paradigms~\cite{chenobject,li2025maniptrans,mandi2025dexmachina}, allowing robots to autonomously explore feasible motion strategies under the guidance of kinematic reference trajectories. These methods commonly design general reward functions encompassing object tracking, hand configurations, and collision dynamics. Maximizing such rewards drives the robot toward successful execution of complex manipulation tasks. Nonetheless, existing works~\cite{mandi2025dexmachina,li2025maniptrans,chenobject} typically draw on limited human motion data sources and some have not transferred the trained robot behaviors to the physical world. Such limitations not only hinder the evaluation of whether the learned policies exhibit behaviorally plausible performance in the real world, but also preclude the integration of sim2real methodologies necessary for robust policy control and for deployment across various environmental conditions. Furthermore, current methods for sim2real transfer of bimanual dexterous manipulation predominantly rely on explicit extraction of object and robot state information~\cite{lin2024learning,chenobject,chen2023sequential,lin2024twisting}, thus failing to achieve end-to-end visual learning. This limitation inherently confines learned policies to specific fixed setups, significantly hindering their adaptability to diverse scenarios.

Motivated by these challenges, we propose \ourshort, a versatile human-to-robot embodied learning framework tailored for mobile bimanual dexterous hand manipulation. \ourshort offers the following three advantages:
\begin{itemize}
    \item \textbf{Diverse sources of human motion}: Our framework supports several human motion sources, including teleoperated simulation data, motion capture~(mocap) data, and raw human videos. We also provide corresponding approaches for data acquisition, enabling \ourshort to efficiently transform varied human motion data into robot-feasible behaviors through RL. Furthermore, these tasks share a uniform set of reward terms, obviating the necessity of designing intricate and task-specific reward functions. In contrast to the methods that depend on collecting a large amount of demonstrations,  we can achieve generalizable policy by augmenting a single reference human motion trajectory coupling with RL training.
    \item \textbf{End-to-end vision-based sim2real transfer}: \ourshort facilitates robust vision-based sim2real transfer by employing DAgger distillation, which converts state-based expert policies into vision-based student policies. Moreover, we introduce a generalized, object-centric depth image augmentation and hybrid control approach, effectively bridging the perception and dynamic sim2real gap.
    \item  \textbf{Mobile manipulation capability}: Our method endows robots with mobile manipulation skills. Building upon ViNT~\cite{shah2023vint}, we develop a RGB-D based module for precise localization wherein the task is modeled as a Perspective-n-Point~(PnP) problem and addressed through an iterative process. This ensures seamless integration with subsequent manipulation tasks and unlock the policy’s capacity to operate autonomously across a broad spectrum of real-world environments.
    
\end{itemize}

With the integration of these capabilities, \ourshort is empowered to execute a wide range of complex or long-horizon mobile bimanual dexterous manipulation tasks across varied and unstructured real-world environments. We conduct extensive experiments in both simulation and real-world settings. The experimental results demonstrate that \ourshort achieves superior performance with high task success rates and high sample efficiency. The learned policies can not only successfully transfer to the physical world but also exhibit generalization capabilities. Furthermore, the navigation component of \ourshort demonstrates precise localization ability, facilitating effective deployment of trained manipulation policies in diverse in-the-wild scenarios.
